\subsection{Autenticação e Autorização (RBAC)}
\label{sec:autenticacao}

A segurança de acesso é definida no \textit{backend} e refletida no cliente. No servidor, o fluxo de autenticação JWT compreende: (i) \textbf{registro}, com a senha \textit{hasheada} com Bcrypt (10 \textit{salt rounds}) antes da persistência; (ii) \textbf{login}, que valida credenciais e emite um \textit{access token} de curta duração (1 hora) e um \textit{refresh token} de longa duração (7 dias); (iii) \textbf{refresh}, em que o cliente troca o \textit{refresh token} por um novo par; e (iv) \textbf{proteção de rotas}, na qual o \texttt{JwtAuthGuard} valida o \textit{access token} e responde 401 quando inválido ou expirado.

O pipeline de segurança é composto por quatro \textit{guards} executados em sequência (Código~\ref{cod:auth_pipeline}):

\begin{lstlisting}[caption={Pipeline de autorização do backend}, label={cod:auth_pipeline}]
Request
  -> JwtAuthGuard      # valida JWT; constroi TenantContext -> req.context
  -> RolesGuard        # verifica @Roles(); isDev ignora
  -> TenantFilterGuard # :organizationId param == req.context.organizationId
  -> DevGuard          # apenas isDev=true passa em rotas @Dev()
  -> Controller
\end{lstlisting}

O \texttt{JwtAuthGuard} enriquece o \textit{payload} JWT no momento do login com todos os dados necessários (\texttt{userId}, \texttt{organizationId}, \texttt{role}, \texttt{isDev}), eliminando consultas ao banco a cada requisição. Usuários listados em \texttt{DEV\_EMAILS} recebem \texttt{isDev=true} no token e ignoram validações de tenant e role. Em termos de papéis, o \textbf{ADMIN} tem acesso total (organizações, usuários, motoristas, configurações), o \textbf{DRIVER} tem acesso limitado (visualizar viagens e confirmar ações relacionadas) e o \textbf{USER} implícito (passageiro) pode inscrever-se em viagens e visualizar suas inscrições. As \textit{memberships} permitem múltiplas \textit{roles} por usuário em diferentes organizações (ex.: motorista em uma org, admin em outra).

\textbf{Reflexo no cliente.} A aplicação persiste a sessão em \texttt{localStorage} (\textit{token} de acesso, \textit{token} de renovação e dados básicos do usuário) e provê o estado de autenticação por meio de dois contextos encadeados: um de autenticação, que envolve um de papel, garantindo que a determinação de papel só ocorra quando a identidade já está disponível. A principal característica de resiliência é a \textbf{renovação transparente de sessão}: diante de uma resposta \texttt{401}, o cliente dispara uma única renovação --- \textbf{deduplicada} entre requisições concorrentes --- e repete a chamada original, conforme o Código~\ref{cod:fe_api}.

\begin{lstlisting}[language=JavaScript, caption={Cliente HTTP com renovação transparente de sessão (\texttt{lib/api.ts})}, label={cod:fe_api}]
export async function api(path, init = {}) {
  const { auth = true, _retry = false, headers, ...rest } = init;
  if (auth && tokenStorage.access)
    headers.Authorization = `Bearer ${tokenStorage.access}`;

  const res = await fetch(`${API_BASE_URL}${path}`, { ...rest, headers });

  // Tenta uma unica renovacao (deduplicada) em caso de 401 e repete a chamada.
  if (res.status === 401 && auth && !_retry) {
    if (!refreshPromise)
      refreshPromise = doRefresh().finally(() => { refreshPromise = null; });
    try { await refreshPromise; return api(path, { ...init, _retry: true }); }
    catch { /* renovacao falhou - propaga o erro original */ }
  }
  if (!res.ok) throw new ApiError(message, res.status, data, errorCode);
  return data;
}
\end{lstlisting}

